\documentclass[10pt,a4paper]{amsart}
\usepackage[margin=19mm]{geometry}
\usepackage{amsmath,amssymb,mathtools}
\usepackage[scr=rsfs,cal=euler]{mathalfa}
\usepackage{xfrac}
\usepackage[unicode,hidelinks,bookmarks=false]{hyperref}
\newcommand{\ie}{i.e.\ }
\newcommand{\cf}{cf.\ }
\newcommand{\resp}{respectively\ }
\newcommand{\Subsection}{\subsection}
\providecommand{\nobreakdash}{\nobreak\hbox{-}}
\setlength{\emergencystretch}{2em}
\multlinegap0pt
\usepackage{bm}
\usepackage{bbm}
\usepackage{dsfont}
\usepackage{xspace}
\usepackage{cancel}
\usepackage{mathtools}
\usepackage{amsthm}
\makeatletter
\@ifundefined{theorem}{\newtheorem{theorem}{Theorem}}{}
\@ifundefined{lemma}{\newtheorem{lemma}[theorem]{Lemma}}{}
\@ifundefined{proposition}{\newtheorem{proposition}[theorem]{Proposition}}{}
\makeatother
\newcommand{\HH}{\mathcal{H}}
\newcommand{\PP}{\mathcal{P}}
\newcommand{\Z}{\mathbb{Z}}
\newcommand{\su}{\subseteq}
\title[A Sharp Ramsey Theorem for Ordered Hypergraph Matchings]{A Sharp Ramsey Theorem for Ordered Hypergraph Matchings}
\author{Lisa Sauermann, Dmitrii Zakharov}
\date{Source: arXiv:2309.04813v4 (CC BY 4.0)}
\begin{document}
\maketitle

\noindent\emph{Source and licence.} A Sharp Ramsey Theorem for Ordered Hypergraph Matchings. Version arXiv:2309.04813v4 (CC BY 4.0), distributed under the Creative Commons Attribution 4.0 International licence, as recorded in the arXiv licence metadata for this exact version and shown on the abstract page https://arxiv.org/abs/2309.04813v4. Full rights notice: licenses/arxiv-2309.04813-CC-BY-4.0.txt.\par\smallskip

\noindent\textit{Editorial scope.} Counted unit: the Theorem whose source label is \texttt{thm-induction-upper-bound} in the article (source0.tex lines 419--425, source label \texttt{thm-induction-upper-bound}), delivered together with the article's own complete original proof of it (source0.tex lines 429--445, one \texttt{proof} environment of 17 lines). No mathematics is written, repaired or re-derived: statement and proof are the article's own text line for line. The proof is detached from the statement in the source: the article states the result at lines 419--425 and prints its proof at lines 429--445, and the proof environment's own header names the statement, so the association is the article's own. Required context retained uncounted: obs-interval-wise (lines 359--362); prop-two-options (lines 381--387). Every definition, display and label of those lines is retained unchanged. Omitted from this package and not redistributed: 1--358, 363--380, 388--418, 426--428, 446--724. Nothing inside a retained excerpt is omitted or altered: every retained body is delivered exactly as the article's lines are written, so no omission receipt of any kind accompanies this package. Citations: the retained excerpts carry no citation command at all, so no bibliography entry of the article is transcribed and no reference is redistributed. Reference adaptation: none was needed and none was made; every reference inside the delivered unit resolves inside it, so no unresolved reference or citation is produced. Printed slips kept literal: no printed word, symbol, hypothesis or number of the counted statement or of its proof was altered. Layout and class: the article's own class is \texttt{aic} with packages including amsthm, amsfonts, amssymb, amsmath, graphicx, tikz, babel, fontenc, inputenc; the wrapper is the lane's pinned \texttt{amsart} editorial wrapper at 10pt on A4, which restates the theorem-like environments the retained text needs and adds the article's own macro definitions that the retained text needs (\texttt{\textbackslash HH}, \texttt{\textbackslash PP}, \texttt{\textbackslash Z}, \texttt{\textbackslash su}), with extra wrapper packages (bm, bbm, dsfont, xspace, cancel, mathtools). The delivered numbering is the unit's own. Assets: no retained span contains a figure environment, a graphics-inclusion command, a PDF-page inclusion, a table or a TikZ picture; no image file is copied into this package, the package declares no asset dependency and no image file is redistributed with it. Licence: version arXiv:2309.04813v4 of this article is distributed under the Creative Commons Attribution 4.0 International licence, as recorded in the arXiv licence metadata for this exact version and shown on the abstract page https://arxiv.org/abs/2309.04813v4. A full rights notice is delivered at licenses/arxiv-2309.04813-CC-BY-4.0.txt. Characters: every retained excerpt is pure ASCII, so the delivered engine, XeLaTeX with its default Latin Modern text font, reports no missing character.
\begin{lemma}\label{obs-interval-wise}
Let $r$ be a positive integer, and let $\lambda=(\lambda_1,\dots,\lambda_s)$ be an ordered partition of $r$. Furthermore, consider a positive integer $m_P$ for every $r$-pattern $P\in \PP(\lambda)$. Let $\HH$ be a finite collection of pairwise disjoint subsets $e\su \Z$, each of size $|e|=r$, such that $\HH$ is interval-wise $\lambda$-partite. For every $P\in \PP(\lambda)$, assume that every $P$-clique contained in $\HH$ has size at most $m_{P}$. Then we have
\[|\HH|\le \prod_{P\in \PP(\lambda)} m_{P}.\]
\end{lemma}

\begin{proposition}\label{prop-two-options}
Let $r$ be a positive integer, and let $\lambda=(\lambda_1,\dots,\lambda_s)$ be an ordered partition of $r$ into $s<r$ parts. For every ordered partition $\lambda'$ of $r$ into $s+1$ parts with $\lambda\succ\lambda'$, consider some positive integer $M_{\lambda'}$, and let $M=\sum_{\lambda'} M_{\lambda'}$ be the sum of all these positive integers. Let $\HH$ be a finite collection of pairwise disjoint subsets $e\su \Z$, each of size $|e|=r$, such that $\HH$ is $\lambda$-partite. Then at least one of the following two statements holds:
\begin{itemize}
    \item[(a)] There is a sub-collection $\HH^*\su \HH$ of size $|\HH^*|\ge |\HH|/(2M)$ such that $\HH^*$ is interval-wise $\lambda$-partite.
    \item[(b)] For some ordered partition $\lambda'$ of $r$ into $s+1$ parts with $\lambda\succ\lambda'$, there is a sub-collection $\HH'\su \HH$ of size $|\HH'|>M_{\lambda'}$ such that $\HH'$ is $\lambda'$-partite.
\end{itemize}
\end{proposition}

\begin{theorem} \label{thm-induction-upper-bound}
Let $r$ be a positive integer, and let $\lambda=(\lambda_1,\dots,\lambda_s)$ be an ordered partition of $r$. For every collectable $r$-pattern $P$, consider a positive integer $m_P$. Let $\HH$ be a finite collection of pairwise disjoint subsets $e\su \Z$, each of size $|e|=r$, such that $\HH$ is $\lambda$-partite. Furthermore, assume that for every collectable $r$-pattern $P$, every $P$-clique contained in $\HH$ has size at most $m_P$. Then the size of $\HH$ is bounded by
\begin{equation}
    |\HH|\le 2^{r-s} \sum_{\lambda^{(s)} \succ \lambda^{(s+1)} \succ \dots \succ \lambda^{(r)}}\ \prod_{ P\in \PP(\lambda^{(s)})\cup\dots\cup \PP(\lambda^{(r)}) } m_{P},
\end{equation}
where the sum is over all sequences $\lambda^{(s)} \succ \lambda^{(s+1)} \succ \dots \succ \lambda^{(r)}$ of ordered partitions of $r$ with $\lambda^{(s)}=\lambda$ (note that then for $t=s,\dots,r$, the ordered partition $\lambda^{(t)}$ must automatically have $t$ parts).
\end{theorem}

\begin{proof}[Proof of Theorem~\ref{thm-induction-upper-bound}]
    We prove the theorem by induction on $r-s$ (note that we always have $s\le r$). If $s=r$, i.e.\ if $\lambda=(1,1,\dots,1)$ is the unique ordered partition of $r$ into $r$ parts, then $\HH$ is automatically interval-wise $\lambda$-partite. Thus, by Lemma~\ref{obs-interval-wise} we have
    \[|\HH|\le \prod_{P\in \PP(\lambda)} m_{P}=2^{r-r} \sum_{\lambda^{(s)} \succ \dots \succ \lambda^{(r)}}\  \prod_{P\in \PP(\lambda^{(s)})\cup\dots\cup \PP(\lambda^{(r)})} m_{P},\]
    where the sum has only one summand corresponding to the unique length-1 sequence $\lambda^{(s)} \succ \dots \succ \lambda^{(r)}$ with $\lambda^{(s)}=\lambda^{(r)}=\lambda$.

    So let us now assume that $s<r$ and that we already proved the desired statement for all smaller values of $r-s$.  For every ordered partition $\lambda'$ of $r$ into $s+1$ parts with $\lambda\succ\lambda'$, let us define
    \[M_{\lambda'}=2^{r-s-1} \sum_{\lambda^{(s+1)} \succ \lambda^{(s+2)} \succ \dots \succ \lambda^{(r)}}\ \prod_{P\in \PP(\lambda^{(s+1)})\cup\dots\cup \PP(\lambda^{(r)})} m_{P},\]
    where the sum is over all sequences $\lambda^{(s+1)} \succ \lambda^{(s+2)} \succ \dots \succ \lambda^{(r)}$ of ordered partitions of $r$ with $\lambda^{(s+1)}=\lambda'$. As in Proposition~\ref{prop-two-options}, let $M=\sum_{\lambda'}M_{\lambda'}$ be the sum of these $M_{\lambda'}$ for all ordered partitions $\lambda'$ of $r$ into $s+1$ parts with $\lambda\succ\lambda'$. Then we have
    \[M=\sum_{\lambda'} M_{\lambda'}=2^{r-s-1} \sum_{\lambda^{(s)}\succ \lambda^{(s+1)}\succ \dots \succ \lambda^{(r)}}\ \prod_{ P\in \PP(\lambda^{(s+1)})\cup\dots\cup \PP(\lambda^{(r)}) } m_{P},\]
    where on the right-hand side the sum is over all sequences $\lambda^{(s)} \succ \lambda^{(s+1)} \succ \dots \succ \lambda^{(r)}$ of ordered partitions of $r$ with $\lambda^{(s)}=\lambda$.

    For every ordered partition $\lambda'$ of $r$ into $s+1$ parts with $\lambda\succ\lambda'$, we can apply the inductive assumption for $\lambda'$, and conclude that for every $\lambda'$-partite sub-collection $\HH'\su \HH$, we must have $|\HH'|\le M_{\lambda'}$. Thus, when applying Proposition~\ref{prop-two-options} to $\HH$ and $\lambda$, with the numbers $M_{\lambda'}$ as defined above, option (b) cannot happen. Hence option (a) must happen, and we can find a sub-collection $\HH^*\su \HH$ of size $|\HH^*|\ge |\HH|/(2M)$ such that $\HH^*$ is interval-wise $\lambda$-partite. But now by Lemma~\ref{obs-interval-wise}, we have
    \[|\HH|/(2M)\le |\HH^*|\le \prod_{P\in \PP(\lambda)} m_{P}.\]
    This gives
    \[|\HH|\le 2M\cdot \prod_{P\in \PP(\lambda)} m_{P}=2^{r-s} \sum_{\lambda^{(s)} \succ \lambda^{(s+1)} \succ \dots \succ \lambda^{(r)}}\ \prod_{ P\in \PP(\lambda^{(s)})\cup\dots\cup \PP(\lambda^{(r)}) } m_{P},\]
    as desired, where the sum on the right-hand side is again over all sequences $\lambda^{(s)} \succ \lambda^{(s+1)} \succ \dots \succ \lambda^{(r)}$ of ordered partitions of $r$ with $\lambda^{(s)}=\lambda$.
    \end{proof}

\end{document}
